\section{Limitations and Next Steps}
\label{sec:limitations}

The local seed pilot demonstrates that the fixed-object coding, validation, and analysis pipeline can be executed and can expose diagnostic failures; it doesn't validate a broader interpretation or use of the assessment. It uses 18 hand-authored seed items, three local models, harmless toy payloads, and one expert adjudicator, the author, labelling items he wrote from expected-behaviour fields he also wrote. The judge pass shares this circularity: the judge model was drawn from the evaluated set and saw each item's expected-behaviour field.

Study A likewise evaluates one realized output for each fixed prompt--model cell; it doesn't estimate generation-level instability across repeated samples or added non-evidential context. Study B is a separate pre-collection design for controlled causal contrasts and can't repair those historical limits retroactively. The results reported here establish that the development pipeline executes, preserve bookkeeping for the observed development contrasts, and motivate judge validation. They don't establish frontier-model safety differences, contextual robustness, contrast stability, or final benchmark labels.

Study B's next empirical pass is a 12-base construction pilot, with four conditions per base, run against several public or accessible models. It should compare three scoring routes on a common response set: expert evaluation, LLM-as-judge annotation, and rule-aided scoring from item metadata. Further construction is governed by the public construct-domain matrix rather than a target item count. A production freeze should fill defensible gaps across item, phenomenon family, source relation, pragmatic status, model, application surface, criterion, and scoring procedure; a larger count doesn't repair missing content.

The later 24-base development set should include independent policy and linguistic evaluation for each item, ideally with two evaluations in each criterion family. It shouldn't collect ordinary-user labels or treat evaluator responses as data about evaluator populations. Ten to twenty percent of items should be repeated or lightly paraphrased to measure item and protocol stability. Released benchmark labels should be produced only after the disagreement source has been classified, with panel and adjudication provenance retained.

Any future two-item seed replacement, outside Study B's four-condition sets, also needs a stricter pair validator. It should require exactly two items per pair, matching phenomenon labels, a declared control contrast, and at least one expected label that changes across the pair.

The current seed P008 is excluded from joint pair completion scoring because it
doesn't meet those conditions; the later benchmark expansion should replace it with two
clean contrast pairs.

AP-SEED-009 needs a different repair. Its three responses split on the antecedent rather than converging, so the disagreement is evidence that the stipulated clarification reference may not be the only defensible one. Study~A re-adjudicates the item output-blind, deciding whether \mention{the previous instruction} is underdetermined enough to require clarification or whether a recency reading is licensed, and its \texttt{capability\_failure} attribution is re-opened at the same time. Revising the expected behaviour now, after the responses are known, would create a new claim rather than repair the old one.

Every planned projection is declared separately in the machine-readable register. The reference bearer is the licensed-action relation generated by a frozen stipulated-regime template and output-blind adjudication procedure. The behavioural bearer is a versioned evaluation-configuration family with enumerated wrappers and context-assembly rules. The measurement bearer is a versioned scorer, codebook, implementation, adjudication procedure, criterion semantics, and enumerated prompt and evidence-access perturbations. The taxonomic bearer is a versioned taxonomy and assignment procedure. Each claim separately names its source result, projected outcome, unit, target cases, transformations, time/version, tolerance, uses, minimum useful reach, evidence, threats, failure triggers, retest, and retirement rule. Post hoc narrowing after target outcomes creates a new claim; it doesn't vindicate the failed one.

LLM-judge validation is an experiment, not a shortcut. The three-by-two pass supplies the disjoint-model and no-rubric conditions, and it shows that a more capable disjoint judge can beat the self-grading one while still discarding the minority class. The next pass needs comparison against independently adjudicated labels rather than provisional author labels, more than one judge prompt per cell, and enough items for the rubric effect to be estimated at better than 18 clusters. It should also test whether a judge that never emits a middle category can be made to use it, since aggregate agreement rewards collapsing the distinction. It should also include a rubric-departure condition, a rubric deliberately misaligned with frontier safety priors, since judges largely keep applying trained priors when the stated definition departs from them \parencite{alloula2026safetyContextual}. It should perturb source-order presentation, include fluent but instruction-violating outputs, and compare judge labels with adjudicated expert labels. Report judge-prompt sensitivity and the failure types that fool the judge, especially quotation/use errors, source-authority errors, information-flow errors under metalinguistic refusal, and over-refusal of safe analysis requests.

Study B construction should counterbalance harmless tokens, source names, ordering, and syntactic forms; reserve untouched templates and payload families; and keep the operative string identical byte for byte where the contrast permits. The lineage manifest binds content hashes to construction splits, and the stored schedule freezes token, source, position, syntax, and marker assignments. Nuisance paraphrases and matched placebos shouldn't change the reference. Blocking shortcut probes cover lexical substitution, position, formatting, source name, length, output vocabulary, operativity-marker removal, and operativity-marker reversal. The pair manifest declares which transformations change the reference and which have to preserve it. The strongest result is a reference-concordant response change exceeding matched nuisance and placebo changes on held-out cases. It supports selective behavioural sensitivity under a named configuration, not a claim that a model internally recognizes authority or possesses a general pragmatic competence.

The minimum Study B development table should report raw outcome vectors and paired transitions by phenomenon family, item, and named configuration, with uncertainty rather than only aggregate correctness. A second table should report nuisance and placebo shifts, including harmful tails. A third should report adjudication instability by item family and criterion conflict. A fourth should identify cases where the adjudicated reference or taxonomy assignment shifted, because those are reference- or taxonomy-projectibility failures rather than model errors.

A separate application-surface table should report prompt-only items against webpage, document, email, tool-result, and transcript wrappers, with risk types split into integrity, confidentiality, tool misuse, policy bypass, and evaluator deception. Surface movement is a direct target test, not automatic transfer. This prevents colour-token success from being mistaken for agent-security robustness. The construction pilot should also report truncation, failed calls, all prespecified repeats, and the shortcut-audit conditions; selectively repeating only failing cells would change the estimand.

The claim register fixes quantitative and categorical failure gates before any production outcome is viewed. One gate illustrates the form: in every claimed family-by-surface cell, at least three of four independent bases must show a behavioural margin whose simultaneous 95\% interval lower bound exceeds 0.20, where that margin subtracts the larger matched nuisance or placebo shift from the observed \(C0-C1\) shift, and each arm must separately reach main-reference correctness of 0.90. A fall in the old token isn't enough on its own. The gates are conjunctive rather than components of a pooled score, so opposite effects can't cancel and equal failures don't count as invariance. The supplement states the full set: the complete base gate, the blocking shortcut probes, the measurement and taxonomic thresholds, the six primary summaries, the manifest-binding conditions, and the \texttt{NOT\_ESTIMATED} rule that applies because the committed record contains no target observations.

Those numbers need their status stated, because this paper argues elsewhere that thresholds encode utilities and intended uses rather than falling out of a response vector. They're provisional preregistration choices, not derived minima. They were set high enough that a passing cell would be worth acting on, and frozen before any target outcome was viewed, which is what stops them being tuned to results. They aren't yet tied to a stated decision context or to a measured cost of a missed block against an over-block, and no sensitivity analysis has been run. The construction pilot should report each cell at nearby thresholds so a reader can see which conclusions depend on the particular values, and any threshold that turns out to carry a conclusion on its own should be re-derived from a declared use and minimum useful reach rather than kept because it was declared first.

Neither the seed pilot nor the current static Study~B fixtures support a claim about long-horizon agent behaviour. A later trajectory harness would have to vary and record sequence length, opportunities to retry, child-session and tool affordances, cumulative state, monitor visibility and intervention, and pause, operational-continuation, normative-release, restart, and stopping rules. OpenAI's first-party incident report usefully motivates those fields, but it names no model or checkpoint and supplies no raw trajectories, incident denominator, monitor thresholds, false-positive rates, or quantitative replay results. OpenAI notes that replay rollouts weren't guaranteed to pursue the same action because of randomness and imperfect environment reconstruction \citep{openai2026longHorizonSafety}. The source supports construct design, not prevalence, comparative performance, or safeguard efficacy.
